\clearpage\fancyhead[L]{\small JUNO.}\hypertarget{native-365-lower}{}\label{native-365-lower}
\begin{center}\textbf{JUNO.}\par\vspace{4pt}\hrule\vspace{4pt}\fontsize{10}{12}\selectfont Monatliche Correspondenz. Band XXVI. Seite 297\;..\;299. September 1812.\par\vspace{4pt}\hrule\vspace{4pt}\end{center}

\begin{flushright}Göttingen, 9. Sept. 1812.\end{flushright}
\par\indent Heute möchte ich meine \emph{Juno}-Beobachtungen nachreichen. Diesen kleinen Planeten zu finden, machte dieses Jahr außerordentlich viel Mühe, da er kaum die 11. Größe erreichte und gerade in der Milchstraße stand. Hinzu kam, dass Mondschein und schlechtes Wetter die Suche unterbrachen. So habe ich nur vier Beobachtungen erhalten, die ich aber alle für sehr gut halte.

\begin{center}\fontsize{10}{12}\selectfont\setlength{\tabcolsep}{4pt}\renewcommand{\arraystretch}{1.15}\begin{tabular}{l|c|c|c}
1812 & Mittlere Zeit & Rektasz. & Deklination\\\hline
Juni 7 & $12^{\mathrm u}\;21^{\mathrm m}\;0^{\mathrm s}$ & $282^\circ\;21'\;24''5$ & $4^\circ\;41'\;44''6$ südl.\\
13 & $12\;51\;12$ & $281\;16\;32.9$ & $4\;33\;39.6$\\
Juli 12 & $11\;35\;28$ & $275\;5\;10.3$ & $5\;6\;11.2$\\
13 & $11\;28\;12$ & $274\;52\;49.9$ & $5\;9\;20.0$\\
\end{tabular}\end{center}

\Needspace{13\baselineskip}\par\indent Meine X Elemente ergeben folgende Unterschiede:

\begin{center}\fontsize{10}{12}\selectfont\setlength{\tabcolsep}{7pt}\renewcommand{\arraystretch}{1.15}\begin{tabular}{l|c|c}
& Rektasz. & Dekl.\\\hline
Juni 7 & $+21''8$ & $-35''9$\\
13 & $+31.0$ & $-34.7$\\
Juli 12 & $+22\;4$ & $-46.9$\\
13 & $+26.2$ & $-45.0$\\
\end{tabular}\end{center}

\par\indent Herr \textsc{Wachter} übernahm es, aus diesen Beobachtungen die Opposition zu berechnen. Ich empfahl ihm dafür eine Methode, die ich seit Langem zu diesem Zweck verwende. Sie hat den Vorteil, dass sich auch zeitlich etwas weiter entfernte Beobachtungen ohne Bedenken für die Opposition verwenden lassen. Ich werde Ihnen ein anderes Mal etwas ausführlicher über dieses Verfahren schreiben. Das Ergebnis von Herrn \textsc{Wachter} lautet:

\Needspace{8\baselineskip}\begin{center}
\emph{Opposition der Juno} 1812.
\end{center}

den 29. Juni $19^{\mathrm u}\;51^{\mathrm m}\;43^{\mathrm s}$ Mittlere Zeit in Göttingen

\begin{center}\fontsize{10}{12}\selectfont\setlength{\tabcolsep}{2pt}\renewcommand{\arraystretch}{1.15}\begin{tabular}{lcr}
Wahre Länge & $.\;.\;.\;.\;.\;.\;.\;.\;.\;.\;.$ & $278^\circ\;15'\;21''3$\\
Wahre geozentrische Breite & $.\;.\;.\;.\;.\;.$ & $18\;36\;56.9$\\
\end{tabular}\end{center}

\Needspace{19\baselineskip}\par\indent Herr \textsc{Wachter} hat außerdem aus den vier letzten beobachteten Oppositionen (die von 1810 wurde bekanntlich versäumt) die Elemente nach der im I. Band der \emph{Comment. Nov. Soc. Sc. Götting.} [S. 1 dieses Bandes] dargestellten Methode korrigiert und Folgendes gefunden:
\SourceLeaf{366}\begin{center}\fontsize{9}{11}\selectfont\setlength{\tabcolsep}{2pt}\renewcommand{\arraystretch}{1.15}\begin{tabular}{lcr}
Epoche 1813, Göttinger Meridian & $.\;.\;.\;.\;.\;.\;.$ & $342^\circ\;56'\;32''2$\\
Tägliche mittlere tropische Bewegung & $.\;.\;.\;.\;.\;.$ & $813''25748$\\
Perihel 1813 & $.\;.\;.\;.\;.\;.\;.\;.\;.\;.\;.\;.\;.\;.$ & $53^\circ\;16'\;52''3$\\
Knoten 1813 & $.\;.\;.\;.\;.\;.\;.\;.\;.\;.\;.\;.\;.\;.\;.\;.$ & $171\;10\;58.9$\\
Neigung der Bahn & $.\;.\;.\;.\;.\;.\;.\;.\;.\;.\;.\;.\;.$ & $13\;4\;17.2$\\
Exzentrizitätswinkel & $.\;.\;.\;.\;.\;.\;.\;.\;.\;.\;.\;.$ & $14\;43\;58.1$\\
Logarithmus der großen Halbachse & $.\;.\;.\;.\;.\;.\;.$ & $0.4265679$\\
\end{tabular}\end{center}

\par\indent Da jetzt die säkularen Änderungen der Elemente der Ceres- und Pallasbahn bekannt sind, soweit sie von Jupiter herrühren (gegen dessen Wirkung die übrigen Planeten ganz unbedeutend sind), veranlasste ich Herrn \textsc{Enke}, einen talentvollen jungen Mann und ebenso geschickten wie sorgfältigen Rechner, zu untersuchen, ob der Abstand der Ceres- und Pallasbahn, der bekanntlich am aufsteigenden Knoten der Ceresbahn auf der Pallasbahn ziemlich klein ist — eben das brachte Herrn Dr. \textsc{Olbers} auf seine bekannte Hypothese —, jetzt ab- oder zunimmt, das heißt, ob die Bahnen sich auf einen wirklichen Schnitt zubewegen oder sich von ihm entfernen. Letzteres würde offenkundig \textsc{Olbers}’ Hypothese stützen, aber das Ergebnis ist genau umgekehrt. Der Abstand war früher größer als heute. Herr \textsc{Enke} findet folgende Radiusvektoren

\begin{center}\fontsize{9}{11}\selectfont\setlength{\tabcolsep}{2pt}\renewcommand{\arraystretch}{1.15}\begin{tabular}{c|c|c|c||c|c|c}
& \multicolumn{3}{c||}{\shortstack{im $\AscendingNode$ der Ceresbahn auf der\\Pallasbahn}} & \multicolumn{3}{c}{\shortstack{im $\DescendingNode$ der Ceresbahn auf der\\Pallasbahn}}\\
Jahr & Ceres & Pallas & Unterschied & Ceres & Pallas & Unterschied\\\hline
$808$ & $2.82294$ & $2.70322$ & $-0.11972$ & $2.67612$ & $2.37780$ & $-0.29832$\\
$1808$ & $2.92427$ & $2.84945$ & $-0.07482$ & $2.59569$ & $2.40346$ & $-0.19223$\\
$3475$ & $2.95374$ & $2.95743$ & $+0.00369$ & $2.57987$ & $2.49163$ & $-0.08824$\\
\end{tabular}\end{center}

\par\indent Danach würde etwa um das Jahr 3397 ein wirklicher Schnitt im $\AscendingNode$ eintreten, was man immerhin als näherungsweise richtig ansehen darf. Freilich wird es auch schon einmal einen Schnitt gegeben haben. Aus dem Verlauf der Zahlen in der vierten und sechsten Spalte lässt sich jedoch wenigstens schließen, dass dies nur zu einer viele Jahrtausende entfernten Epoche möglich gewesen sein kann. Wenn man also die Hypothese von Dr. \textsc{Olbers} über den Ursprung der neuen Planeten annehmen will, liegt dieser Ursprung in einer noch gar nicht berechenbaren Ferne vor der Zeit, bis zu der unsere Geschichte zurückreicht.
